\section*{Sets and Indices}

\begin{itemize}
    \item $I$: set of food items from \texttt{table\_1.csv}, indexed by $i$.
    \item $V \subseteq I$: subset of vegetable items, i.e., items with \texttt{Type = Vegetable}.
\end{itemize}

\section*{Parameters}

For each item $i \in I$:
\begin{itemize}
    \item $p_i$: protein content per 100g (\texttt{Protein\_Content\_per\_100g}).
    \item $c_i$: cost per 100g in USD (\texttt{Cost\_per\_100g}).
    \item $k_i$: calories per 100g in kcal (\texttt{Calories\_per\_100g}).
    \item $s_i$: sodium per 100g in mg (\texttt{Sodium\_mg\_per\_100g}).
\end{itemize}

Global parameters from \texttt{general\_parameters.csv}:
\begin{itemize}
    \item $B$: maximum daily budget (\texttt{max\_budget}).
    \item $W$: total daily purchase weight limit in grams (\texttt{total\_weight\_limit}).
    \item $m^{\mathrm{veg}}$: minimum number of vegetable types required in each meal (\texttt{min\_vegetable\_types}).
    \item $K^{\min}$: minimum daily calories from consumed food (\texttt{min\_daily\_calories}).
    \item $K^{\max}$: maximum daily calories from consumed food (\texttt{max\_daily\_calories}).
    \item $S^{\max}$: maximum daily sodium from consumed food (\texttt{max\_daily\_sodium}).
\end{itemize}

Derived constants used exactly as in code:
\begin{itemize}
    \item $U = \left\lfloor \dfrac{W}{100} \right\rfloor$ (\texttt{max\_units}): item-wise upper bound on daily purchased 100g units.
    \item $M = 3$ (\texttt{max\_meal\_units}): item-wise upper bound on lunch or dinner 100g units.
\end{itemize}

All unit variables below are measured in 100g units.

\section*{Decision Variables}

For each item $i \in I$:
\begin{itemize}
    \item $q_i$ (code: \texttt{purchase\_units[i]}) $\in \mathbb{Z}_{\ge 0}$: total purchased units of item $i$ for the day.
    \item $\ell_i$ (code: \texttt{lunch\_units[i]}) $\in \mathbb{Z}_{\ge 0}$: units of item $i$ assigned to lunch.
    \item $d_i$ (code: \texttt{dinner\_units[i]}) $\in \mathbb{Z}_{\ge 0}$: units of item $i$ assigned to dinner.
    \item $z_i$ (code: \texttt{select\_item[i]}) $\in \{0,1\}$: equals 1 if item $i$ is selected for purchase.
\end{itemize}

For each vegetable item $i \in V$:
\begin{itemize}
    \item $y_i^L$ (code: \texttt{lunch\_veg\_selected[i]}) $\in \{0,1\}$: equals 1 if vegetable $i$ appears in lunch.
    \item $y_i^D$ (code: \texttt{dinner\_veg\_selected[i]}) $\in \{0,1\}$: equals 1 if vegetable $i$ appears in dinner.
\end{itemize}

\section*{Objective}

Maximize total consumed protein over lunch and dinner:
\[
\max \sum_{i \in I} p_i(\ell_i + d_i).
\]

\section*{Constraints}

Daily budget:
\[
\sum_{i \in I} c_i q_i \le B.
\]

Daily purchase weight limit:
\[
100 \sum_{i \in I} q_i \le W.
\]

Purchase/use linking for each item $i \in I$:
\[
q_i \ge \ell_i + d_i,
\]
\[
q_i \le U z_i,
\]
\[
\ell_i \le M z_i,
\]
\[
d_i \le M z_i.
\]

Exact meal sizes:
\[
\sum_{i \in I} \ell_i = 3,
\]
\[
\sum_{i \in I} d_i = 3.
\]

Daily calorie bounds, computed from consumed units only:
\[
\sum_{i \in I} k_i(\ell_i + d_i) \ge K^{\min},
\]
\[
\sum_{i \in I} k_i(\ell_i + d_i) \le K^{\max}.
\]

Daily sodium bound, computed from consumed units only:
\[
\sum_{i \in I} s_i(\ell_i + d_i) \le S^{\max}.
\]

Vegetable-selection linking for each $i \in V$:
\[
\ell_i \ge y_i^L,
\]
\[
\ell_i \le M y_i^L,
\]
\[
d_i \ge y_i^D,
\]
\[
d_i \le M y_i^D.
\]

Minimum number of vegetable types in each meal:
\[
\sum_{i \in V} y_i^L \ge m^{\mathrm{veg}},
\]
\[
\sum_{i \in V} y_i^D \ge m^{\mathrm{veg}}.
\]

Integrality and binary restrictions:
\[
q_i,\ell_i,d_i \in \mathbb{Z}_{\ge 0} \qquad \forall i \in I,
\]
\[
z_i \in \{0,1\} \qquad \forall i \in I,
\]
\[
y_i^L,y_i^D \in \{0,1\} \qquad \forall i \in V.
\]

\section*{Notes on Implementation}

\begin{itemize}
    \item The code formulates a mixed-integer linear program and solves it with \texttt{GUROBI\_CMD(msg=0)} through the PuLP-compatible interface.
    \item Meal portions are enforced in 100g units as $\sum_i \ell_i = 3$ and $\sum_i d_i = 3$, which corresponds to exactly 300g for lunch and 300g for dinner.
    \item The variable $q_i$ may exceed $\ell_i+d_i$, so the implementation allows purchased-but-unused food.
    \item The binary purchase variable $z_i$ is linked by upper bounds only; together with nonnegativity and $q_i \ge \ell_i+d_i$, this ensures that if $z_i=0$ then $q_i=\ell_i=d_i=0$.
    \item For vegetable items, the binaries $y_i^L$ and $y_i^D$ are linked both above and below to meal usage, so for integer $\ell_i,d_i$ they indicate whether the corresponding vegetable is used positively in lunch or dinner.
    \item The model has no explicit requirement that a meal contain any protein item; it only enforces meal size, vegetable-count minima, daily nutrition bounds, budget, and purchase-linking constraints exactly as implemented.
\end{itemize}
